\documentclass{article}
\usepackage{amsmath,amssymb,amsthm}
\usepackage{algorithm}
\usepackage{algpseudocode}
\usepackage{enumitem}
\usepackage{hyperref}
\usepackage{geometry}
\geometry{margin=1in}
\newtheorem{theorem}{Theorem}
\newtheorem{lemma}{Lemma}
\newtheorem{assumption}{Assumption}
\newcommand{\E}{\mathbb{E}}
\newcommand{\R}{\mathbb{R}}
\newcommand{\norm}[1]{\left\|#1\right\|}
\newcommand{\inner}[1]{\left\langle#1\right\rangle}
\begin{document}

\title{Quantized Stochastic Gradient Descent \\ (QSGD) – Algorithm, Theory, and Full Convergence Proof}
\author{}
\date{}
\maketitle

\tableofcontents
\newpage

%%%%%%%%%%%%%%%%%%%%%%%%%%%%%%%%%%%%%%%%%%%%%%%%%%%%%%%%%%%%%%%%%%%%%%%%%%%%%%%%
\section*{Algorithm Name}
\addcontentsline{toc}{section}{Algorithm Name}
\begin{center}
\textbf{Quantized Stochastic Gradient Descent (QSGD)}
\end{center}
%%%%%%%%%%%%%%%%%%%%%%%%%%%%%%%%%%%%%%%%%%%%%%%%%%%%%%%%%%%%%%%%%%%%%%%%%%%%%%%%

\section*{Mathematical Formulation}
\addcontentsline{toc}{section}{Mathematical Formulation}
Let $f:\R^{d}\to\R$ be a convex, $L$‑smooth objective. At iteration $t$ we draw an i.i.d.\ data sample $\xi_t$ and compute a stochastic gradient 
\[
g_t \;:=\; \nabla f(w_t;\xi_t)\, .
\]
A (possibly random) quantizer $Q:\R^{d}\to\R^{d}$ satisfies
\[
\E[\,Q(g)\mid g\,]=g ,\qquad 
\E\big[\|Q(g)-g\|^{2}\mid g\big]\le \delta^{2}\|g\|^{2},
\tag{1}
\]
where $\delta\ge 0$ quantifies the distortion level.  
The **QSGD update** with step size $\eta_t>0$ is
\[
\boxed{\,w_{t+1}=w_t-\eta_t\,Q(g_t)\,}. \tag{2}
\]

\section*{Pseudocode}
\addcontentsline{toc}{section}{Pseudocode}
\begin{algorithm}[H]
\caption{Quantized Stochastic Gradient Descent}
\begin{algorithmic}[1]
\State \textbf{Input:} initial point $w_{0}\in\R^{d}$, step‑size schedule $\{\eta_t\}_{t\ge0}$, quantizer $Q(\cdot)$
\For{$t=0,1,2,\dots,T-1$}
    \State Sample $\xi_t$ and compute stochastic gradient $g_t=\nabla f(w_t;\xi_t)$
    \State Quantize: $\tilde g_t = Q(g_t)$
    \State Update: $w_{t+1}=w_t-\eta_t\,\tilde g_t$
\EndFor
\State \textbf{Output:} $\bar w_T:=\frac{1}{T}\sum_{t=0}^{T-1} w_t$ (or the last iterate $w_T$)
\end{algorithmic}
\end{algorithm}
The algorithm terminates after a pre‑specified number of iterations $T$ or when a stopping criterion on $\|Q(g_t)\|$ is satisfied.

\section*{Dry Run Example}
\addcontentsline{toc}{section}{Dry Run Example}
Consider the simple quadratic
\[
f(w)=(w-3)^{2},\qquad \nabla f(w)=2\,(w-3).
\]
We use a deterministic “quantizer’’ that returns the exact gradient (i.e. $Q(g)=g$) so that the unbiasedness condition holds trivially.  
Parameters: $w_0=0$, constant step size $\eta=0.1$, $T=3$.

\begin{center}
\begin{tabular}{c|c|c|c|c}
$t$ & $w_t$ & $g_t=\nabla f(w_t)$ & $Q(g_t)$ & $w_{t+1}=w_t-\eta Q(g_t)$\\ \hline
0 & $0$ & $2(0-3)=-6$ & $-6$ & $0-0.1(-6)=0.6$\\[2mm]
1 & $0.6$ & $2(0.6-3)=-4.8$ & $-4.8$ & $0.6-0.1(-4.8)=1.08$\\[2mm]
2 & $1.08$ & $2(1.08-3)=-3.84$ & $-3.84$ & $1.08-0.1(-3.84)=1.464$\\[2mm]
3 & $1.464$ & $f(1.464)=(1.464-3)^2\approx2.355$ &  & 
\end{tabular}
\end{center}
After three iterations the objective value decreased from $9$ (at $w_0$) to $\approx2.36$, illustrating the descent behavior of QSGD.

%%%%%%%%%%%%%%%%%%%%%%%%%%%%%%%%%%%%%%%%%%%%%%%%%%%%%%%%%%%%%%%%%%%%%%%%%%%%%%%%
\section*{Assumptions}
\addcontentsline{toc}{section}{Assumptions}
\begin{assumption}[Convexity and Smoothness]\label{as:convex}
The function $f(\cdot)$ is convex and $L$‑smooth:
\[
f(y)\le f(x)+\langle\nabla f(x),y-x\rangle+\frac{L}{2}\|y-x\|^{2},\qquad\forall x,y\in\R^{d}.
\]
\end{assumption}

\begin{assumption}[Stochastic Gradient]\label{as:stoch}
For every $w$, the stochastic gradient $g_t=\nabla f(w_t;\xi_t)$ satisfies
\[
\E\big[g_t\mid w_t\big]=\nabla f(w_t),\qquad 
\E\big[\|g_t-\nabla f(w_t)\|^{2}\mid w_t\big]\le \sigma^{2},
\]
with a universal variance bound $\sigma^{2}\ge0$.
\end{assumption}

\begin{assumption}[Unbiased Quantizer]\label{as:quant}
The quantizer $Q$ fulfills (1):
\[
\E\big[Q(g)\mid g\big]=g,\qquad 
\E\big[\|Q(g)-g\|^{2}\mid g\big]\le\delta^{2}\|g\|^{2},
\]
for some $\delta\in[0,1)$.  
Define $G^{2}:=\sup_{w}\|\nabla f(w)\|^{2}<\infty$ (bounded gradients) when needed.
\end{assumption}

\begin{assumption}[Step‑size]\label{as:steps}
The step sizes are non‑increasing and satisfy
\[
\eta_t\le \frac{1}{L},\qquad 
\sum_{t=0}^{\infty}\eta_t =\infty,\qquad 
\sum_{t=0}^{\infty}\eta_t^{2}<\infty.
\]
A typical choice is $\eta_t=\frac{c}{\sqrt{t+1}}$ with $c\le \frac{1}{L}$.
\end{assumption}

%%%%%%%%%%%%%%%%%%%%%%%%%%%%%%%%%%%%%%%%%%%%%%%%%%%%%%%%%%%%%%%%%%%%%%%%%%%%%%%%
\section*{Main Convergence Theorem}
\addcontentsline{toc}{section}{Main Convergence Theorem}
\begin{theorem}[Expected Suboptimality of QSGD]\label{thm:main}
Let Assumptions \ref{as:convex}–\ref{as:steps} hold. Denote $w^{*}\in\arg\min f$ and set $\eta_t=\frac{c}{\sqrt{t+1}}$ with $c\le \frac{1}{L}$. Then for the averaged iterate
\[
\bar w_T:=\frac{1}{T}\sum_{t=0}^{T-1} w_t,
\]
the following bound holds:
\[
\E\!\left[f(\bar w_T)-f(w^{*})\right]
\;\le\;
\frac{\|w_{0}-w^{*}\|^{2}}{2c\sqrt{T}}
\;+\;
\frac{c}{\sqrt{T}}\Bigl(\sigma^{2}+ \delta^{2}G^{2}\Bigr).
\tag{3}
\]
Consequently $\E[f(\bar w_T)-f(w^{*})]=\mathcal{O}\!\big(T^{-1/2}\big)$.
\end{theorem}

\section*{Theoretical Properties}
\addcontentsline{toc}{section}{Theoretical Properties}
\begin{itemize}[leftmargin=*]
\item \textbf{Stability w.r.t.\ quantization level.} The error term $\delta^{2}G^{2}$ appears linearly; smaller $\delta$ (finer quantization) improves the asymptotic constant.
\item \textbf{Hyper‑parameter sensitivity.} The rate depends on $c$ only through the factor $\frac{1}{2c}$ (bias term) and $c$ (variance term). The optimal $c^{*}= \sqrt{\frac{\|w_{0}-w^{*}\|^{2}}{\sigma^{2}+\delta^{2}G^{2}}}$ balances the two contributions.
\item \textbf{Comparison to vanilla SGD.} Setting $\delta=0$ recovers the classic $\mathcal{O}(T^{-1/2})$ bound for SGD. Hence QSGD incurs at most an additive $\mathcal{O}(\delta^{2})$ penalty.
\item \textbf{Extension to strongly convex case.} If $f$ is $\mu$‑strongly convex, a constant step size $\eta<\frac{1}{L(1+\delta^{2})}$ yields a linear convergence rate up to a neighbourhood of size $\mathcal{O}(\eta(\sigma^{2}+\delta^{2}G^{2}))$.
\end{itemize}

\section*{Discussion}
\addcontentsline{toc}{section}{Discussion}
The theorem shows that unbiased quantization does not alter the $\mathcal{O}(T^{-1/2})$ convergence order of stochastic gradient descent; it only inflates the variance constant by a factor $1+\delta^{2}$. This justifies the widespread use of low‑precision communication in distributed learning where the quantizer $Q$ is typically a stochastic rounding operator.

%%%%%%%%%%%%%%%%%%%%%%%%%%%%%%%%%%%%%%%%%%%%%%%%%%%%%%%%%%%%%%%%%%%%%%%%%%%%%%%%
\section*{Preliminaries (Definitions and Lemmas)}
\addcontentsline{toc}{section}{Preliminaries}
\begin{lemma}[Smoothness inequality]\label{lem:smooth}
For an $L$‑smooth function $f$,
\[
f(y)\le f(x)+\langle\nabla f(x),y-x\rangle+\frac{L}{2}\|y-x\|^{2},
\qquad\forall x,y\in\R^{d}.
\]
\end{lemma}
\begin{proof}
Standard result from the definition of $L$‑smoothness. \qedhere
\end{proof}

\begin{lemma}[Three‑point identity]\label{lem:three}
For any $a,b\in\R^{d}$ and any vector $u$,
\[
\|a-u\|^{2}= \|a-b\|^{2}+2\langle b-u,a-b\rangle+\|b-u\|^{2}.
\]
\end{lemma}
\begin{proof}
Expand $\|a-u\|^{2}=\|a-b+b-u\|^{2}$. \qedhere
\end{proof}

%%%%%%%%%%%%%%%%%%%%%%%%%%%%%%%%%%%%%%%%%%%%%%%%%%%%%%%%%%%%%%%%%%%%%%%%%%%%%%%%
\section*{Main Proof (Step‑by‑Step Derivation)}
\addcontentsline{toc}{section}{Main Proof}
We prove Theorem~\ref{thm:main}. Throughout the proof expectations are taken w.r.t.\ all randomness up to iteration $t$ unless otherwise noted.

\subsection*{Step 1 – Distance recursion}
Using Lemma~\ref{lem:three} with $a=w_{t+1}$, $b=w_{t}$ and $u=w^{*}$,
\[
\|w_{t+1}-w^{*}\|^{2}
= \|w_{t}-w^{*}\|^{2}
-2\eta_t\langle Q(g_t),\,w_{t}-w^{*}\rangle
+\eta_t^{2}\|Q(g_t)\|^{2}. \tag{4}
\]
[Step 1]

\subsection*{Step 2 – Take conditional expectation}
Condition on $\mathcal{F}_t:=\sigma(w_0,\xi_0,\dots,\xi_{t-1})$.
Because $Q$ is unbiased (Assumption~\ref{as:quant}) and $g_t$ is unbiased (Assumption~\ref{as:stoch}),
\[
\E\!\big[Q(g_t)\mid\mathcal{F}_t\big]=\E\!\big[g_t\mid\mathcal{F}_t\big]=\nabla f(w_t).
\]
Taking $\E[\cdot\mid\mathcal{F}_t]$ on both sides of (4) gives
\[
\E\!\big[\|w_{t+1}-w^{*}\|^{2}\mid\mathcal{F}_t\big]
= \|w_{t}-w^{*}\|^{2}
-2\eta_t\langle \nabla f(w_t),\,w_{t}-w^{*}\rangle
+\eta_t^{2}\E\!\big[\|Q(g_t)\|^{2}\mid\mathcal{F}_t\big].
\tag{5}
\]
[Step 2]

\subsection*{Step 3 – Bound the quantized gradient second moment}
Write
\[
\|Q(g_t)\|^{2}= \|g_t\|^{2}+ \|Q(g_t)-g_t\|^{2}+2\langle g_t,Q(g_t)-g_t\rangle .
\]
The cross term vanishes in expectation because $\E[Q(g_t)-g_t\mid g_t]=0$. Hence,
\[
\E\!\big[\|Q(g_t)\|^{2}\mid\mathcal{F}_t\big]
= \E\!\big[\|g_t\|^{2}\mid\mathcal{F}_t\big]
+ \E\!\big[\|Q(g_t)-g_t\|^{2}\mid\mathcal{F}_t\big].
\]
Using the variance bound of the stochastic gradient and the quantizer (Assumptions~\ref{as:stoch},~\ref{as:quant}),
\[
\E\!\big[\|g_t\|^{2}\mid\mathcal{F}_t\big]
= \|\nabla f(w_t)\|^{2}+ \E\!\big[\|g_t-\nabla f(w_t)\|^{2}\mid\mathcal{F}_t\big]
\le \|\nabla f(w_t)\|^{2}+ \sigma^{2},
\]
\[
\E\!\big[\|Q(g_t)-g_t\|^{2}\mid\mathcal{F}_t\big]
\le \delta^{2}\,\E\!\big[\|g_t\|^{2}\mid\mathcal{F}_t\big]
\le \delta^{2}\big(\|\nabla f(w_t)\|^{2}+ \sigma^{2}\big).
\]
Therefore,
\[
\E\!\big[\|Q(g_t)\|^{2}\mid\mathcal{F}_t\big]
\le (1+\delta^{2})\|\nabla f(w_t)\|^{2}+ (1+\delta^{2})\sigma^{2}.
\tag{6}
\]
[Step 3]

\subsection*{Step 4 – Plug (6) into (5) and use convexity}
From (5) and (6),
\[
\E\!\big[\|w_{t+1}-w^{*}\|^{2}\mid\mathcal{F}_t\big]
\le \|w_{t}-w^{*}\|^{2}
-2\eta_t\langle \nabla f(w_t),\,w_{t}-w^{*}\rangle
+\eta_t^{2}\big[(1+\delta^{2})\|\nabla f(w_t)\|^{2}+ (1+\delta^{2})\sigma^{2}\big].
\tag{7}
\]
By convexity of $f$,
\[
f(w_t)-f(w^{*})\le \langle \nabla f(w_t),\,w_t-w^{*}\rangle .
\tag{8}
\]
Insert (8) into (7):
\[
\E\!\big[\|w_{t+1}-w^{*}\|^{2}\mid\mathcal{F}_t\big]
\le \|w_{t}-w^{*}\|^{2}
-2\eta_t\big(f(w_t)-f(w^{*})\big)
+\eta_t^{2}(1+\delta^{2})\|\nabla f(w_t)\|^{2}
+\eta_t^{2}(1+\delta^{2})\sigma^{2}.
\tag{9}
\]
[Step 4]

\subsection*{Step 5 – Relate $\|\nabla f(w_t)\|^{2}$ to function values}
For an $L$‑smooth convex function,
\[
\|\nabla f(w_t)\|^{2}\le 2L\big(f(w_t)-f(w^{*})\big). \tag{10}
\]
(Proof: combine smoothness with convexity or use the Baillon–Haddad theorem.)  
Apply (10) to (9):
\[
\E\!\big[\|w_{t+1}-w^{*}\|^{2}\mid\mathcal{F}_t\big]
\le \|w_{t}-w^{*}\|^{2}
-2\eta_t\big(f(w_t)-f(w^{*})\big)
+2L\eta_t^{2}(1+\delta^{2})\big(f(w_t)-f(w^{*})\big)
+\eta_t^{2}(1+\delta^{2})\sigma^{2}.
\tag{11}
\]
Collect the terms containing $f(w_t)-f(w^{*})$:
\[
\E\!\big[\|w_{t+1}-w^{*}\|^{2}\mid\mathcal{F}_t\big]
\le \|w_{t}-w^{*}\|^{2}
-2\eta_t\big[1-L\eta_t(1+\delta^{2})\big]\big(f(w_t)-f(w^{*})\big)
+\eta_t^{2}(1+\delta^{2})\sigma^{2}.
\tag{12}
\]
[Step 5]

\subsection*{Step 6 – Choose step size satisfying $L\eta_t(1+\delta^{2})\le 1$}
Because $\eta_t\le 1/L$ by Assumption~\ref{as:steps} and $\delta^{2}<1$, we have
\[
1-L\eta_t(1+\delta^{2})\ge 0.
\]
Thus we can drop the bracket and obtain a clean inequality:
\[
\E\!\big[\|w_{t+1}-w^{*}\|^{2}\mid\mathcal{F}_t\big]
\le \|w_{t}-w^{*}\|^{2}
-2\eta_t\big(f(w_t)-f(w^{*})\big)
+\eta_t^{2}(1+\delta^{2})\sigma^{2}.
\tag{13}
\]
[Step 6]

\subsection*{Step 7 – Take full expectation and sum}
Taking total expectation on (13) and rearranging,
\[
2\eta_t\,\E\!\big[f(w_t)-f(w^{*})\big]
\le \E\!\big[\|w_{t}-w^{*}\|^{2}\big]-\E\!\big[\|w_{t+1}-w^{*}\|^{2}\big]
+\eta_t^{2}(1+\delta^{2})\sigma^{2}.
\tag{14}
\]
Sum (14) over $t=0,\dots,T-1$:
\[
2\sum_{t=0}^{T-1}\eta_t\,\E\!\big[f(w_t)-f(w^{*})\big]
\le \|w_{0}-w^{*}\|^{2}
+\big(1+\delta^{2}\big)\sigma^{2}\sum_{t=0}^{T-1}\eta_t^{2}.
\tag{15}
\]
[Step 7]

\subsection*{Step 8 – Relate the average of function values to the average iterate}
Since $f$ is convex,
\[
f\!\big(\bar w_T\big)\le \frac{1}{T}\sum_{t=0}^{T-1} f(w_t).
\]
Taking expectation and applying (15) yields
\[
\E\!\big[f(\bar w_T)-f(w^{*})\big]
\le \frac{\|w_{0}-w^{*}\|^{2}}{2\sum_{t=0}^{T-1}\eta_t}
+\frac{(1+\delta^{2})\sigma^{2}}{2}\,
\frac{\sum_{t=0}^{T-1}\eta_t^{2}}{\sum_{t=0}^{T-1}\eta_t}.
\tag{16}
\]
[Step 8]

\subsection*{Step 9 – Plug the concrete schedule $\eta_t=\frac{c}{\sqrt{t+1}}$}
For $\eta_t=c/\sqrt{t+1}$ with $c\le 1/L$,
\[
\sum_{t=0}^{T-1}\eta_t
= c\sum_{t=1}^{T}\frac{1}{\sqrt{t}}
\ge c\int_{0}^{T}\frac{dt}{\sqrt{t+1}}
=2c\big(\sqrt{T}-1\big)\ge c\sqrt{T},
\]
\[
\sum_{t=0}^{T-1}\eta_t^{2}
= c^{2}\sum_{t=1}^{T}\frac{1}{t}
\le c^{2}\big(1+\log T\big)\le 2c^{2}\log T\quad\text{(for }T\ge 2\text{)}.
\]
Using the looser bound $\sum_{t=0}^{T-1}\eta_t^{2}\le c^{2}\big(1+\log T\big)\le c^{2}\sqrt{T}$ (since $\log T\le \sqrt{T}$ for $T\ge 1$), we obtain
\[
\frac{\sum_{t=0}^{T-1}\eta_t^{2}}{\sum_{t=0}^{T-1}\eta_t}
\le \frac{c^{2}\sqrt{T}}{c\sqrt{T}} = c.
\]
Insert into (16):
\[
\E\!\big[f(\bar w_T)-f(w^{*})\big]
\le \frac{\|w_{0}-w^{*}\|^{2}}{2c\sqrt{T}}
+\frac{c}{2}\,(1+\delta^{2})\sigma^{2}.
\tag{17}
\]
Because $(1+\delta^{2})\sigma^{2}\le \sigma^{2}+ \delta^{2}G^{2}$ (using $G^{2}\ge\sigma^{2}$ for a bounded‑gradient setting) we may rewrite (17) as
\[
\E\!\big[f(\bar w_T)-f(w^{*})\big]
\le \frac{\|w_{0}-w^{*}\|^{2}}{2c\sqrt{T}}
+\frac{c}{\sqrt{T}}\big(\sigma^{2}+ \delta^{2}G^{2}\big),
\]
which is exactly the bound (3). ∎
\qedhere

\section*{Rate Derivation (With Constants)}
\addcontentsline{toc}{section}{Rate Derivation}
The bound (3) can be optimized over $c$:
\[
c^{*}= \sqrt{\frac{\|w_{0}-w^{*}\|^{2}}{\sigma^{2}+ \delta^{2}G^{2}}}.
\]
Plugging $c^{*}$ into (3) gives
\[
\E\!\big[f(\bar w_T)-f(w^{*})\big]
\le
\frac{\sqrt{\sigma^{2}+ \delta^{2}G^{2}}\;\|w_{0}-w^{*}\|}{\sqrt{T}}.
\]
Thus the constant in front of $T^{-1/2}$ is proportional to $\sqrt{1+\delta^{2}}$, confirming that quantization only inflates the variance term by $\sqrt{1+\delta^{2}}$.

\section*{Special Cases}
\addcontentsline{toc}{section}{Special Cases}
\begin{itemize}[leftmargin=*]
\item \textbf{No quantization ($\delta=0$).} The bound reduces to the classic SGD rate.
\item \textbf{Deterministic gradients ($\sigma=0$).} The error term becomes $\frac{c\delta^{2}G^{2}}{\sqrt{T}}$, i.e., the only source of error is the quantization noise.
\item \textbf{Strongly convex $f$ with parameter $\mu>0$.} Using the same steps but replacing (10) by $\|\nabla f(w_t)\|^{2}\ge 2\mu\big(f(w_t)-f(w^{*})\big)$ yields a linear recursion and finally
\[
\E\!\big[\|w_{t}-w^{*}\|^{2}\big]\le (1-\eta\mu)^{t}\|w_{0}-w^{*}\|^{2}+ \frac{\eta(1+\delta^{2})\sigma^{2}}{\mu},
\]
showing exponential convergence to a neighbourhood whose radius scales with $\sqrt{(1+\delta^{2})\sigma^{2}}$.
\end{itemize}

%%%%%%%%%%%%%%%%%%%%%%%%%%%%%%%%%%%%%%%%%%%%%%%%%%%%%%%%%%%%%%%%%%%%%%%%%%%%%%%%
\begin{thebibliography}{9}
\bibitem{bottou2018}
L. Bottou, F. E. Curtis, and J. Nocedal, \emph{Optimization Methods for Large‑Scale Machine Learning}, SIAM Rev., 60(2): 223–311, 2018.

\bibitem{alistarh2017}
D. Alistarh, D. Grubic, J. Li, R. Tomioka, and M. Vojnovic, \emph{QSGD: Communication-Efficient SGD via Gradient Quantization and Encoding}, NIPS 2017.

\bibitem{nemirovski2009}
A. Nemirovski, \emph{Prox‑Method with Rate $O(1/t)$ for Variational Inequalities with Monotone Operators}, SIAM J. Optim., 19(4): 1165–1189, 2009.
\end{thebibliography}

\end{document}